\section{Wochen Dokumentation}
\label{sec:wochen}

Die folgende Dokumentation hält den Projektfortschritt Woche für Woche
fest. Jede Woche ist nach demselben Schema aufgebaut: Ziel, durchgeführte
Arbeiten, Ergebnisse, Probleme und nächste Schritte. So lässt sich der
Verlauf des Projekts auch im Nachhinein gut nachvollziehen.

\subsection{KW40 Grundgerüst des FNN}

\paragraph{Ziel der Woche.}
Aufbau eines ersten lauffähigen Feedforward-Netzes mit beliebig vielen
Hidden Layers, zunächst ohne Bias.

\paragraph{Durchgeführte Arbeiten.}
\begin{table}[htbp]
    \centering
    \small
    \begin{tabular}{l p{9cm}}
    \toprule
    \textbf{Person} & \textbf{Aufgabe} \\
    \midrule
    Anton   & Implementierung der Klasse \texttt{FNW} mit
              beliebig vielen Hidden Layers; Forward Pass mit
              Gewichten (ohne Bias) \\
    Aaron   & Recherche zu Aktivierungsfunktionen (ReLU, Sigmoid, Tanh)
              und deren Ableitungen \\
    Liam    & Erstellung des Projektbuch-Gerüsts und erste Dokumentation
              der Theorie \\
    \bottomrule
    \end{tabular}
    \caption{Aufgabenverteilung KW40.}
    \label{tab:kw40}
\end{table}

\paragraph{Ergebnisse.}
\begin{itemize}
    \item Lauffähige Klasse \texttt{FNW} mit konfigurierbarer
          Layer-Struktur.
    \item Forward Pass funktioniert für Netze mit $n$ Hidden Layers.
    \item Theoretische Grundlagen zu Aktivierungsfunktionen dokumentiert.
\end{itemize}


\paragraph{Nächste Schritte.}
\begin{itemize}
    \item Ergänzung des Bias-Terms im Forward Pass.
    \item Vorbereitung der Loss-Berechnung und Backpropagation.
\end{itemize}

\subsection{KW41}

\paragraph{Platzhalter}

\newpage